\documentclass[11pt,letterpaper]{article}
\usepackage[T1]{fontenc}
\usepackage{lmodern,amsmath,amssymb,amsthm,microtype}
\usepackage[margin=0.92in]{geometry}
\usepackage[hidelinks]{hyperref}
\newtheorem{theorem}{Theorem}
\title{Weighted rank pooling for asymmetric diagonal-gated posets}
\author{Aafi Khan Malek\\\small Independent Researcher, Ahmedabad, India\\
\small\href{mailto:aafimalek2032@gmail.com}{aafimalek2032@gmail.com}}
\date{November 2026}
\begin{document}
\maketitle
\begin{abstract}
A checked weighted rank-pooling inequality broadens an earlier
diagonal-gate criterion for the existence of a balanced comparison
in a finite poset. For a prefix of length \(h\) and offsets \(d,e\),
the new sufficient condition is \(de\le(h-1)(d+e)\). No bound on
rail length is needed. We give the counting argument, state the
exact Lean endpoint, and distinguish this class from the unrestricted
width-three conjecture.
\end{abstract}

\section*{Statement}
A comparison is balanced if the probability of its chosen order in a
uniform random linear extension is in \([1/3,2/3]\). Let \(a,b\) be
incomparable in a finite poset. Every strict predecessor of \(a\)
is strictly below \(b\). Assume that the probability of \(a<b\)
is at most \(1/2\). The exclusive branch
\[
 S=\{z:\ b\le z,\ a\not\le z\}
\]
consists of a common \(b\)-first chain prefix of length \(h\ge2\),
below two internally ordered rails \(X=(x_0,\ldots)\) and
\(Y=(y_0,\ldots)\). The two rails partition the remaining branch
points; either may be empty. For nonnegative integers \(d,e\),
assume \(x_i<y_{i+d}\) and \(y_j<x_{j+e}\) whenever the indicated
endpoints exist. Additional cross relations are allowed.

\begin{theorem}
If \(de\le(h-1)(d+e)\), some \(z\in S\) makes the comparison
\(a<z\) balanced. This suffices for arbitrary finite rail lengths.
\end{theorem}
The \href{https://github.com/Aafimalek/novel-discoveries/tree/b6e15ab63380419495d7fe9813427c704283e869/width-three-posets/verification/diagonal-gated-two-arm-branches/lean-asymmetric-weighted-class}{Lean endpoint}
is \texttt{WidthThree.Diagonal.Labels.balanced\_of\_diagonal\_product}.
The independent typed audit states
\(N\le3C(a<z)\le2N\) for ranking count \(N\).
The strict replay on Lean 4.36.0 uses only
\texttt{propext}, \texttt{Classical.choice}, and \texttt{Quot.sound}.
The exact gate and good-pair hypotheses are essential to the stated
scope; general width three is not claimed.

\section*{Weighted count transfer}
Fix one induced branch order and let \(c_k\) count ambient rankings
with exactly \(k\) branch points before \(a\) in that order.
The \href{https://github.com/Aafimalek/novel-discoveries/commit/2c65834be82b6d2f1dc072a26c8f4c7994bcb946}{fixed-order swap lemma}
makes \(c_k\) nonincreasing. Append zero beyond the last cut
and put \(w_t=c_{t-1}-c_t\ge0\). For a branch point of one-based
rank \(r\), the count for \(a\) before that point in the fixed order is
\[
 F(r)=\sum_{k<r}c_k=\sum_{t\ge1}w_t\min(r,t).
\]
Suppose every induced order satisfies
\(p r_x+q r_y\le(p+q)(r_u+r_v)\), for fixed nonnegative
integer weights \(p,q\). For each threshold \(t\),
\[
 p\min(r_x,t)+q\min(r_y,t)
 \le(p+q)\min(r_u+r_v,t)
 \le(p+q)(\min(r_u,t)+\min(r_v,t)).
\]
The first inequality follows by bounding the left side both by
\(p r_x+q r_y\) and by \((p+q)t\). Summing with coefficients \(w_t\)
and then across induced orders gives the corresponding ambient count
inequality. This does not require independent cuts or equal weighting
of distinct induced orders.

\section*{Application to the two first high comparisons}
The \href{https://github.com/Aafimalek/novel-discoveries/tree/5b1476e2b00d787dc6ae630fefd4b8510d0e50e3/width-three-posets/verification/diagonal-gated-two-arm-branches/lean-complete-rank-only-class}{checked
unweighted assembly} handles the prefix, final cuts, and the cases
with no high comparison or with high comparisons in just one rail.
For the remaining case, assume no balanced comparison and take the
first high points \(x_i,y_j\), where high means probability strictly
above \(2/3\) and low means strictly below \(1/3\). Put
\[
 B=h+i+j+1,\quad \delta=i-j,\quad
 M=\max\{\min(i,j),|\delta|\}.
\]
The fixed two low anchors \(u,v\) selected in the
\href{https://github.com/Aafimalek/novel-discoveries/tree/1ec4dfe997e57556815c849c30fb2cc4c384428e/width-three-posets/verification/diagonal-gated-two-arm-branches/lean-all-fixed-anchors}{checked
anchor argument} have \(r_u+r_v\ge B+(h-1)+M\) in every induced order.
The earlier head has rank at most \(B\). If \(x_i\) is later its
rank is at most \(B+d+\delta\); if \(y_j\) is later the bound is
\(B+e-\delta\). These gate bounds remain valid at finite ends:
if a gate endpoint is absent the opposing rail ends before that
index. Weight the first high ranks by \(e,d\), respectively. Their
weighted excess over \((d+e)B\) is at most
\(de+e\delta\) or \(de-d\delta\), in either case at most
\((d+e)((h-1)+M)\). Weighted pooling therefore contradicts
the two high comparisons and the two low anchor comparisons
when \(d+e>0\). If both offsets vanish and there are two heads,
the gates at index zero would force opposite strict orders.
This completes the change to the existing finite case assembly.

\paragraph{Scope and credit.}
When \(h=4\), offsets \((4,12)\) satisfy the product criterion, but
not the old cap \(\max(d,e)\le 2h-2\). More generally, if one
offset is at most \(h-1\), this criterion does not bound the other.
This comparison concerns our specific criteria, not a priority claim
over the broader balanced-pair literature. See
\href{https://arxiv.org/abs/2609.30888}{Aires, Chan, Pak and Panova}
and \href{https://zenodo.org/records/24357683}{Fathi's working paper}
for related results with different hypotheses. The research note,
argument and Lean development used AI assistance.
\end{document}
